\documentclass[11pt,a4paper]{article}
\usepackage[margin=2.5cm]{geometry}
\usepackage[utf8]{inputenc}
\usepackage[T1]{fontenc}
\usepackage{lmodern}
\usepackage{microtype}
\usepackage{amsmath,amssymb,amsthm}
\usepackage{mathtools}
\usepackage{booktabs}
\usepackage{enumitem}
\usepackage[numbers,sort&compress]{natbib}
\usepackage{xcolor}
\usepackage{hyperref}
\hypersetup{
  colorlinks=true,
  linkcolor=blue!50!black,
  citecolor=blue!50!black,
  urlcolor=blue!50!black,
  pdftitle={The Ideal Gas Laws from a Bounded Phase Space Resolved at a Fixed Resolution},
  pdfauthor={Kundai Farai Sachikonye}
}

\newtheorem{theorem}{Theorem}[section]
\newtheorem{proposition}[theorem]{Proposition}
\newtheorem{lemma}[theorem]{Lemma}
\newtheorem{corollary}[theorem]{Corollary}
\theoremstyle{definition}
\newtheorem{definition}[theorem]{Definition}
\newtheorem{axiom}{Axiom}
\renewcommand{\theaxiom}{R\arabic{axiom}}
\newtheorem{hypothesis}{Hypothesis}
\renewcommand{\thehypothesis}{\Alph{hypothesis}}
\theoremstyle{remark}
\newtheorem{remark}[theorem]{Remark}

\newcommand{\R}{\mathbb{R}}
\newcommand{\N}{\mathbb{N}}
\newcommand{\X}{\mathcal{X}}
\newcommand{\Sym}{\mathfrak{S}}
\newcommand{\kB}{k_{\mathrm{B}}}
\newcommand{\dd}{\mathrm{d}}
\newcommand{\Pmech}{P_{\mathrm{mech}}}

\title{\textbf{The Ideal Gas Laws from a Bounded Phase Space\\Resolved at a Fixed Resolution}}
\author{
Kundai Farai Sachikonye\\
AIMe Registry for Artificial Intelligence\\
\texttt{kundai.sachikonye@bitspark.com}
}
\date{}

\begin{document}
\maketitle

\begin{abstract}
\noindent
The textbook derivation of the ideal gas law rests on five ingredients.
The five ingredients are a phase-space cell of fixed size, the division by $N!$ for identical particles, equal weights for equal phase-space measure, the identification of entropy with the logarithm of a count, and an energy shell of positive thickness.
Each ingredient is usually introduced as a postulate or as a convenience.
This paper derives all five from four stated axioms about an observer that resolves a bounded phase space at a fixed resolution, together with one explicitly stated dynamical hypothesis.
The axioms fix the cell size as a unit, force the energy constraint to have positive width, and require identical particles to carry no labels.
Entropy is defined operationally as the minimum average number of binary distinctions needed to locate the state among the resolved alternatives, and it is shown to equal $\kB\ln K$ to within one distinction, where $K$ is the count of resolved states.
With temperature and pressure defined as the conjugates of energy and volume, the relation $PV=N\kB T$ follows exactly, for every particle number $N\ge1$, every mass, and every energy function that does not depend on position.
The independence of the law from the nature of the particles yields Avogadro's law.
Equipartition holds up to a correction that is bounded explicitly and vanishes as $N^{-1}$.
The mechanical pressure, computed as the time-averaged impulse on the walls, satisfies $\Pmech V=\tfrac23 E$ exactly and agrees with the thermodynamic pressure within the same bound.
For a single particle the mechanical law holds exactly, and its identification with the temperature requires the other particles of the gas.
Dalton's law, the Sackur--Tetrode entropy, the resolution of the Gibbs mixing paradox, the adiabatic law, and the Maxwell--Boltzmann distribution follow within the same framework.
\end{abstract}

\tableofcontents

%======================================================================
\section{Introduction}
\label{sec:intro}
%======================================================================

\subsection{The ingredients of the standard derivation}

The ideal gas law $PV=N\kB T$ relates the pressure $P$, the volume $V$, the particle number $N$ and the temperature $T$ of a dilute gas.
Here $\kB$ is a positive constant that fixes the unit of temperature.
In the statistical derivation due to Boltzmann and Gibbs, and in its modern textbook form \citep{boltzmann1877,gibbs1902,huang1987}, the law is obtained from five ingredients:
\begin{enumerate}[label=(I\arabic*),leftmargin=3em]
\item\label{ing:cell} a phase-space cell of fixed measure $h^{3}$ per particle, which turns a measure into a count;
\item\label{ing:nfact} a division of that count by $N!$ for identical particles;
\item\label{ing:equal} equal weights for regions of equal phase-space measure;
\item\label{ing:log} the identification of the entropy with $\kB$ times the logarithm of the count;
\item\label{ing:shell} an energy shell of positive thickness in place of the energy surface.
\end{enumerate}
Ingredient~\ref{ing:cell} is usually justified by quantum mechanics, \ref{ing:nfact} by indistinguishability or by the requirement that the entropy be extensive, \ref{ing:equal} by an appeal to ignorance or to ergodicity, \ref{ing:log} by Boltzmann's principle, and \ref{ing:shell} by convenience \citep{huang1987,khinchin1949,jaynes1965}.

\subsection{Approach}

This paper treats the five ingredients as consequences rather than as postulates.
The starting point is an observer that resolves a bounded phase space at a fixed, finite resolution.
Four axioms state what such an observer can and cannot name (Section~\ref{sec:setting}).
One dynamical hypothesis, ergodicity on the energy surface, is stated separately because it is the one assumption whose status is not settled for the systems of interest (Hypothesis~\ref{hyp:ergodic}).
From the axioms and the hypothesis, ingredient~\ref{ing:cell} follows as the definition of a unit, \ref{ing:shell} from the positivity of the resolution, and \ref{ing:nfact} from the absence of particle labels.
Ingredient~\ref{ing:equal} follows from the dynamics together with the absence of distinguished cells, and \ref{ing:log} from an operational definition of entropy as a locating cost.

\subsection{Summary of results}

The principal results are the following.
\begin{enumerate}[label=(\roman*),leftmargin=2.2em]
\item The description of an energy at a fixed resolution has positive width (Lemma~\ref{lem:noshell}).
\item The states of identical particles are permutation orbits, and the measure of a symmetric region on the space of orbits is its phase-space measure divided by $N!$ (Lemma~\ref{lem:quotient}).
\item The minimum average number of binary distinctions needed to locate one of $K$ equally weighted states lies in $[\log_2K,\log_2K+1)$ (Theorem~\ref{thm:locating}).
\item With temperature and pressure defined as conjugates, $PV=N\kB T$ holds exactly for every $N\ge1$ (Theorem~\ref{thm:igl}), for every position-independent energy function (Theorem~\ref{thm:carrier}).
\item The energy satisfies $E=\tfrac32N\kB T(1-\eta)$ with an explicit bound on $\eta$ of order $N^{-1}$ (Theorem~\ref{thm:equipartition}).
\item The time-averaged impulse on the walls gives $\Pmech V=\tfrac23E$ exactly for every trajectory without interparticle collisions (Theorem~\ref{thm:mech}).
\item One particle obeys $\Pmech V=\tfrac23\varepsilon$ exactly; within a gas of $N$ particles its time-averaged pressure obeys $p_1V=\kB T$ up to the same correction (Theorems~\ref{thm:single-isolated} and~\ref{thm:single-gas}).
\end{enumerate}

\subsection{Organisation}

Section~\ref{sec:setting} defines the phase space, the ideal gas and the axioms.
Section~\ref{sec:resolution} draws the consequences of a fixed resolution for the energy constraint.
Section~\ref{sec:labels} constructs the states of identical particles and the count of a region.
Section~\ref{sec:weights} obtains the weights of states from the dynamics.
Section~\ref{sec:entropy} defines the entropy as a locating cost and computes it.
Section~\ref{sec:laws} derives the equation of state, equipartition and the adiabatic law.
Section~\ref{sec:mech} computes the mechanical pressure.
Section~\ref{sec:single} treats a single particle and the Maxwell--Boltzmann distribution.
Section~\ref{sec:mixtures} treats mixtures.
Section~\ref{sec:expansion} treats free expansion.
Section~\ref{sec:discussion} states what was assumed and what was derived.

%======================================================================
\section{Setting and axioms}
\label{sec:setting}
%======================================================================

\subsection{Phase space}

\begin{definition}[Container, phase space and Liouville measure]
\label{def:phase}
The \emph{container} is the closed cube $\Lambda=[0,L]^{3}$ of side $L>0$ and volume $V=L^{3}$.
The \emph{single-particle phase space} is $\X_1=\Lambda\times\R^{3}$, whose points $y=(q,p)$ consist of a position $q\in\Lambda$ and a momentum $p\in\R^{3}$.
For an integer $N\ge1$, the \emph{$N$-particle phase space} is the product $\X_N=\X_1^{N}$, whose points are $x=(y_1,\dots,y_N)$ with $y_i=(q_i,p_i)$.
The \emph{Liouville measure} $\mu_N$ is Lebesgue measure on $\X_N\subset\R^{6N}$.
\end{definition}

\begin{definition}[Ideal gas]
\label{def:ideal}
An \emph{ideal gas} of $N$ particles of mass $m>0$ is the phase space $\X_N$ together with the energy function
\begin{equation}
H_N(x)=\sum_{i=1}^{N}\varepsilon(p_i),\qquad \varepsilon(p)=\frac{|p|^{2}}{2m},
\label{eq:hamiltonian}
\end{equation}
and the motion in which each particle moves freely in the interior of $\Lambda$ and reflects specularly at its faces.
Specular reflection at a face reverses the momentum component normal to that face and leaves the other two components unchanged.
The quantity $\varepsilon(p_i)$ is the \emph{kinetic energy} of particle $i$, and $H_N$ is the \emph{energy}.
\end{definition}

The energy~\eqref{eq:hamiltonian} does not depend on the positions.
This independence is the defining property of ideality used below, and Theorem~\ref{thm:carrier} shows that it is the only property of $H_N$ on which the equation of state depends.

\begin{definition}[Energy surface, ball and shell]
\label{def:shell}
For $E>0$, the \emph{energy surface} is $\Sigma_E=\{x\in\X_N:H_N(x)=E\}$ and the \emph{energy ball} is $B_E=\{x\in\X_N:H_N(x)\le E\}$.
For $0<\Delta E\le E$, the \emph{energy shell} of thickness $\Delta E$ is
\begin{equation}
W(E,\Delta E)=\{x\in\X_N:\ E-\Delta E<H_N(x)\le E\}=B_E\setminus B_{E-\Delta E}.
\label{eq:shell}
\end{equation}
Its \emph{relative thickness} is $\delta=\Delta E/E\in(0,1]$.
When $V$ and $N$ must be shown, the shell is written $W(E,\Delta E;V,N)$.
\end{definition}

The volume of the unit ball in $\R^{d}$ is denoted
\begin{equation}
\omega_d=\frac{\pi^{d/2}}{\Gamma\!\left(\tfrac d2+1\right)},
\label{eq:omega}
\end{equation}
where $\Gamma$ is Euler's gamma function, $\Gamma(z)=\int_0^\infty t^{z-1}e^{-t}\,\dd t$ for $z>0$.

\subsection{The observer}

The \emph{observer} is whatever draws distinctions between states of the gas: an instrument, a record, or an agent.
An \emph{observable} is a measurable function $f\colon\X_N\to\R$ whose values the observer can register.
A \emph{description} is a set of states that the observer can name, such as ``the energy lies in a given interval''.
The four axioms constrain what the observer can name.

\begin{axiom}[Boundedness]
\label{ax:bounded}
Every description names a region of $\X_N$ of finite positive Liouville measure.
\end{axiom}

\begin{axiom}[Fixed resolution]
\label{ax:resolution}
The observer resolves $\X_1$ into \emph{cells}, which are measurable sets of equal Liouville measure $h^{3}$ for a fixed constant $h>0$.
The $N$-particle cells are the products of $N$ single-particle cells, of measure $h^{3N}$.
Every description is a union of cells, up to sets of measure zero, and the \emph{count} of a description is its measure in units of the cell measure.
\end{axiom}

\begin{axiom}[No distinguished particle]
\label{ax:labels}
The particles are identical: every observable is invariant under every permutation of the particle indices.
That is, $f(y_{\sigma(1)},\dots,y_{\sigma(N)})=f(y_1,\dots,y_N)$ for every observable $f$ and every permutation $\sigma$ of $\{1,\dots,N\}$.
\end{axiom}

\begin{axiom}[No distinguished cell]
\label{ax:cells}
Within a description, no cell receives more weight than another unless some observable distinguishes them.
\end{axiom}

\begin{remark}[On the axioms]
\label{rem:axioms}
Axiom~\ref{ax:resolution} fixes $h$ as a unit and gives it no further interpretation; the results of Section~\ref{sec:laws} do not depend on its value (Proposition~\ref{prop:hdep}).
Axiom~\ref{ax:labels} states that the index $i$ is a bookkeeping device of the theorist and not a property of the gas.
Assigning a particle a name in advance of any observation that distinguishes it would be an act of selection that the gas does not support.
Axiom~\ref{ax:cells} applies the same principle to cells instead of particles.
\end{remark}

\subsection{The dynamics}

The motion of Definition~\ref{def:ideal} defines a flow $\phi_t\colon\X_N\to\X_N$, $t\in\R$, away from a set of measure zero of trajectories that meet an edge or a corner of $\Lambda$.
The flow conserves $H_N$, so it maps each energy surface $\Sigma_E$ into itself.
It also preserves the Liouville measure.
For free motion this preservation is Liouville's theorem \citep{arnold1989}, and for motion with specular reflection at the faces of a polyhedron it is a standard property of billiard flows \citep{chernov2006}.
The flow therefore preserves the \emph{surface measure} $\nu_E$ defined, for measurable $A\subseteq\X_N$, by
\begin{equation}
\nu_E(A)=\int_{A\cap\Sigma_E}\frac{\dd\sigma}{|\nabla H_N|},
\label{eq:surface-measure}
\end{equation}
where $\sigma$ is the $(6N-1)$-dimensional surface measure on $\Sigma_E$ and $\nabla H_N$ is the gradient of $H_N$ in $\R^{6N}$ \citep{khinchin1949}.

\begin{hypothesis}[Ergodicity on the energy surface]
\label{hyp:ergodic}
For each $E>0$, the flow restricted to $\Sigma_E$ is ergodic with respect to $\nu_E$: every measurable subset of $\Sigma_E$ that the flow maps into itself has $\nu_E$-measure zero or full $\nu_E$-measure.
\end{hypothesis}

\begin{remark}[Status of the hypothesis]
\label{rem:ergodic-status}
Hypothesis~\ref{hyp:ergodic} fails for the strictly non-interacting gas of Definition~\ref{def:ideal}.
Under specular reflection in a cube, the magnitudes $|p_{i,x}|$, $|p_{i,y}|$ and $|p_{i,z}|$ of every momentum component of every particle are conserved separately.
An ideal gas is therefore understood, here as elsewhere, as the limit of a gas of rarely interacting particles.
The interaction is rare enough to make a negligible contribution to the energy and to the pressure, and frequent enough to make the motion ergodic.
For hard disks in a torus, ergodicity has been proved for typical systems \citep{simanyi2003}.
Every result below that uses Hypothesis~\ref{hyp:ergodic} says so.
The results of Sections~\ref{sec:resolution}, \ref{sec:labels}, \ref{sec:entropy}, \ref{sec:laws} and~\ref{sec:mech} do not use it.
\end{remark}

%======================================================================
\section{Consequences of a fixed resolution}
\label{sec:resolution}
%======================================================================

\begin{lemma}[No sharp energy surface]
\label{lem:noshell}
Under Axioms~\ref{ax:bounded} and~\ref{ax:resolution}, no description equals the energy surface $\Sigma_E$ up to a set of measure zero.
Every description that contains $\Sigma_E$ up to a set of measure zero has Liouville measure at least $h^{3N}$.
\end{lemma}

\begin{proof}
The surface $\Sigma_E$ is the product of $\Lambda^{N}$ with a sphere of radius $\sqrt{2mE}$ in $\R^{3N}$, so $\mu_N(\Sigma_E)=0$.
By Axiom~\ref{ax:resolution} a description is a union of cells, and a non-empty union of cells has measure at least $h^{3N}>0$.
No set of positive measure equals a set of measure zero up to a null set.
A description containing $\Sigma_E$ is non-empty, so its measure is at least $h^{3N}$.
\end{proof}

The observer therefore cannot describe the gas as having energy exactly $E$.
The smallest description consistent with energy $E$ is a band of energies of positive width, and this band is modelled here by the shell $W(E,\Delta E)$ of Definition~\ref{def:shell}.
The thickness $\Delta E>0$ is fixed by the resolution and is treated as a parameter.
Theorem~\ref{thm:equipartition} and Corollary~\ref{cor:shell-insensitive} show that the results do not depend on it at leading order in $N$.

\begin{lemma}[Measure of the shell]
\label{lem:shell-measure}
Write $a=3N/2$.
The energy ball and the energy shell have Liouville measures
\begin{align}
\mu_N(B_E)&=V^{N}\,\omega_{3N}\,(2mE)^{a},\label{eq:ball-measure}\\
\mu_N\bigl(W(E,\Delta E)\bigr)&=V^{N}\,\omega_{3N}\,(2mE)^{a}\,\bigl[1-(1-\delta)^{a}\bigr].\label{eq:shell-measure}
\end{align}
Both are finite and positive, so the shell satisfies Axiom~\ref{ax:bounded}.
\end{lemma}

\begin{proof}
The energy does not depend on positions, so $B_E=\Lambda^{N}\times\{p\in\R^{3N}:|p|^{2}\le2mE\}$.
The first factor has measure $V^{N}$ and the second is a ball of radius $\sqrt{2mE}$ in $\R^{3N}$, of measure $\omega_{3N}(2mE)^{3N/2}$.
The shell is $B_E\setminus B_{E-\Delta E}$ with $B_{E-\Delta E}\subseteq B_E$, and $(E-\Delta E)^{a}=E^{a}(1-\delta)^{a}$.
\end{proof}

%======================================================================
\section{States without labels}
\label{sec:labels}
%======================================================================

\subsection{Permutation orbits}

The symmetric group $\Sym_N$ of permutations of $\{1,\dots,N\}$ acts on $\X_N$ by $\sigma\cdot(y_1,\dots,y_N)=(y_{\sigma(1)},\dots,y_{\sigma(N)})$.
The \emph{orbit} of $x$ is $\Sym_N\cdot x=\{\sigma\cdot x:\sigma\in\Sym_N\}$.
A region $A\subseteq\X_N$ is \emph{symmetric} if $\sigma\cdot A=A$ for every $\sigma\in\Sym_N$.
The shells and balls of Definition~\ref{def:shell} are symmetric.

\begin{proposition}[States are orbits]
\label{prop:orbits}
Under Axiom~\ref{ax:labels}, two points of the same orbit take the same value on every observable.
Hence the observer cannot separate them, and its states are the orbits of $\Sym_N$ in $\X_N$.
\end{proposition}

\begin{proof}
If $x'=\sigma\cdot x$, then $f(x')=f(x)$ for every observable $f$ by Axiom~\ref{ax:labels}.
\end{proof}

\begin{lemma}[Measure on the space of orbits]
\label{lem:quotient}
Let $F=\{x\in\X_N:\ q_{1,1}<q_{2,1}<\dots<q_{N,1}\}$, where $q_{i,1}$ is the first Cartesian coordinate of the position of particle $i$.
Then:
\begin{enumerate}[label=(\roman*),leftmargin=2.2em]
\item the sets $\sigma\cdot F$, $\sigma\in\Sym_N$, are pairwise disjoint, and their union is $\X_N$ minus a set of measure zero;
\item every orbit outside that null set meets $F$ in exactly one point;
\item for every symmetric measurable region $A$, $\mu_N(A\cap F)=\mu_N(A)/N!$.
\end{enumerate}
\end{lemma}

\begin{proof}
Let $D=\{x:\ q_{i,1}=q_{j,1}\text{ for some }i\ne j\}$.
Each set $\{q_{i,1}=q_{j,1}\}$ is a hyperplane of $\R^{6N}$, so $\mu_N(D)=0$.
For $x\notin D$ the first coordinates are distinct, and exactly one permutation $\sigma$ sorts them into increasing order, which places $\sigma^{-1}\cdot x$ in $F$.
The uniqueness of the sorting permutation gives (i) and (ii).
For (iii), each permutation preserves $\mu_N$ and $A$ is symmetric, so $\mu_N(A\cap\sigma\cdot F)=\mu_N(\sigma\cdot(A\cap F))=\mu_N(A\cap F)$.
Summing over the $N!$ disjoint sets $\sigma\cdot F$ gives $\mu_N(A)=N!\,\mu_N(A\cap F)$.
\end{proof}

The region $F$ represents the space of orbits: each state of the observer corresponds to one point of $F$, up to measure zero.
The measure of a symmetric region, counted in states rather than in labelled points, is therefore $\mu_N(A)/N!$.

\subsection{The count of a region}

\begin{definition}[Count]
\label{def:count}
The \emph{count} of a symmetric region $A\subseteq\X_N$ of finite measure at resolution $h$ is
\begin{equation}
K(A)=\frac{\mu_N(A)}{N!\,h^{3N}}.
\label{eq:count}
\end{equation}
For the shell, write $K(E,V,N)=K\bigl(W(E,\Delta E;V,N)\bigr)$, with $\Delta E$ held fixed.
\end{definition}

The count~\eqref{eq:count} is the measure of the region on the space of orbits (Lemma~\ref{lem:quotient}), expressed in units of the cell measure (Axiom~\ref{ax:resolution}).
Ingredients~\ref{ing:cell} and~\ref{ing:nfact} of Section~\ref{sec:intro} are thereby obtained from Axioms~\ref{ax:resolution} and~\ref{ax:labels}.

\subsection{Occupation counting in the dilute regime}

The count~\eqref{eq:count} uses measure on the space of orbits.
A second way to count states without labels uses occupations of position cells, and the two agree when the gas is dilute.

\begin{proposition}[Occupation counts]
\label{prop:occupation}
Partition $\Lambda$ into $M$ congruent position cells of volume $v=V/M$, and let $N\le M/2$.
A positional configuration of $N$ unlabelled particles at this resolution is an \emph{occupation function} $n\colon\{1,\dots,M\}\to\{0,1,2,\dots\}$ with $\sum_jn(j)=N$.
Let $C_{\le1}$ be the number of occupation functions with every $n(j)\le1$, and $C$ the number of all occupation functions.
Then
\begin{equation}
\frac{(M-N)^{N}}{N!}\le C_{\le1}=\binom{M}{N}\le\frac{M^{N}}{N!}\le C=\binom{M+N-1}{N}\le\frac{(M+N)^{N}}{N!},
\label{eq:occupation-bounds}
\end{equation}
and both $\ln C_{\le1}$ and $\ln C$ differ from $\ln\bigl(M^{N}/N!\bigr)$ by at most $2N^{2}/M$.
\end{proposition}

\begin{proof}
An occupation function with values in $\{0,1\}$ is a choice of $N$ of the $M$ cells, so $C_{\le1}=\binom MN=\prod_{j=0}^{N-1}(M-j)/N!$.
Each factor lies in $[M-N,M]$, which gives the first two inequalities.
An arbitrary occupation function is a multiset of size $N$ drawn from $M$ cells, so $C=\binom{M+N-1}{N}=\prod_{j=0}^{N-1}(M+j)/N!$.
Each factor lies in $[M,M+N]$, which gives the last two inequalities.
For the logarithms, $0\le-N\ln(1-N/M)\le2N^{2}/M$ when $N/M\le\tfrac12$, since $-\ln(1-u)\le u/(1-u)\le2u$ for $0\le u\le\tfrac12$.
Likewise $0\le N\ln(1+N/M)\le N^{2}/M$.
\end{proof}

The quantity $M^{N}/N!=V^{N}/(N!\,v^{N})$ is the positional factor of~\eqref{eq:count} when the single-particle cell of measure $h^{3}$ is the product of a position cell of volume $v$ and a momentum cell of volume $h^{3}/v$.
Proposition~\ref{prop:occupation} shows that the counts with and without multiple occupation of a cell agree with this factor to within $2N^{2}/M$ in the logarithm.
Per particle the discrepancy is at most $2N/M=2(N/V)v$, the number density times twice the cell volume.
The ideal gas is the regime in which this number is small, and in that regime the statistics of multiple occupation do not affect the count.

%======================================================================
\section{Weights from the traversal}
\label{sec:weights}
%======================================================================

\begin{theorem}[Weights on the energy surface]
\label{thm:weights}
Assume Hypothesis~\ref{hyp:ergodic}.
For every $E>0$, every measurable $A\subseteq\Sigma_E$, and $\nu_E$-almost every initial point $x\in\Sigma_E$,
\begin{equation}
\lim_{\tau\to\infty}\frac1\tau\int_0^{\tau}\mathbf 1_A\bigl(\phi_t(x)\bigr)\,\dd t=\frac{\nu_E(A)}{\nu_E(\Sigma_E)},
\label{eq:birkhoff}
\end{equation}
where $\mathbf 1_A$ is the indicator function of $A$.
More generally, for every $\nu_E$-integrable function $g$ on $\Sigma_E$, the time average of $g$ along almost every trajectory equals the $\nu_E$-average of $g$.
\end{theorem}

\begin{proof}
The flow preserves the finite measure $\nu_E$ and is ergodic with respect to it.
The statement is the ergodic theorem for flows \citep{birkhoff1931,walters1982}.
\end{proof}

\begin{proposition}[Weights on the shell]
\label{prop:shell-weights}
The normalised Liouville measure $\mu_N(\,\cdot\cap W)/\mu_N(W)$ on the shell $W=W(E,\Delta E)$ is invariant under the flow.
It gives equal weight to cells of equal measure, and it decomposes as
\begin{equation}
\mu_N(A\cap W)=\int_{E-\Delta E}^{E}\nu_{E'}(A\cap\Sigma_{E'})\,\dd E'
\label{eq:coarea}
\end{equation}
for every measurable $A$.
\end{proposition}

\begin{proof}
Invariance holds because the flow preserves $\mu_N$ and maps $W$ into itself.
Equal weight for equal measure holds by construction.
Equation~\eqref{eq:coarea} is the coarea formula applied to $H_N$ with the density~\eqref{eq:surface-measure}.
\end{proof}

Theorem~\ref{thm:weights} fixes the weights along each energy surface from the dynamics alone.
The observer cannot resolve which surface of the shell the gas occupies (Lemma~\ref{lem:noshell}).
By Axiom~\ref{ax:cells} it then gives all cells of the shell equal weight, which is the normalised Liouville measure of Proposition~\ref{prop:shell-weights}.
By~\eqref{eq:coarea}, this measure is the superposition, with uniform density in energy, of the surface measures that the dynamics produces.
The two sources of weight are therefore consistent, and ingredient~\ref{ing:equal} is obtained.

\begin{remark}[Visiting does not change the count]
\label{rem:three-counts}
Three different numbers are associated with the gas, and they should not be confused.
The first is the count $K$ of resolved states in a description, which depends only on the description and the resolution.
The second is $\log_2K$, the number of binary distinctions needed to locate one state among them (Section~\ref{sec:entropy}).
The third is the number of times the trajectory enters a given cell during an interval of time, which grows without bound by Poincar\'e recurrence \citep{poincare1890}.
The trajectory enters cells without any act of the observer, and the number of such entries does not alter $K$.
No quantity defined from $K$ changes with the duration of observation.
\end{remark}

%======================================================================
\section{Entropy as a locating cost}
\label{sec:entropy}
%======================================================================

\subsection{Locating a state}

\begin{definition}[Locating procedure and cost]
\label{def:locating}
Let a finite set of $K\ge1$ alternatives be given.
A \emph{locating procedure} is a rooted binary tree whose $K$ leaves are labelled bijectively by the alternatives.
Each internal node is a \emph{binary distinction}: a question with two answers that sends the alternatives consistent with that node to its two children.
The \emph{cost} of locating an alternative is the depth of its leaf, the number of distinctions on the path from the root to it.
\end{definition}

\begin{theorem}[Cost of locating]
\label{thm:locating}
Let $K\ge1$ alternatives carry equal weights $1/K$.
\begin{enumerate}[label=(\roman*),leftmargin=2.2em]
\item Every locating procedure has average cost at least $\log_2K$.
\item There is a locating procedure whose cost is at most $\lceil\log_2K\rceil$ for every alternative, so that the minimum average cost lies in $[\log_2K,\log_2K+1)$.
\end{enumerate}
\end{theorem}

\begin{proof}
(i) Let $d_1,\dots,d_K$ be the leaf depths of a locating procedure.
The leaves of a binary tree satisfy Kraft's inequality $\sum_{j}2^{-d_j}\le1$ \citep{kraft1949,cover2006}.
The function $u\mapsto2^{-u}$ is convex, so by Jensen's inequality $2^{-\bar d}\le\frac1K\sum_j2^{-d_j}\le\frac1K$, where $\bar d=\frac1K\sum_jd_j$.
Hence $\bar d\ge\log_2K$.
(ii) Assign distinct binary strings of length $\lceil\log_2K\rceil$ to the alternatives, which is possible because $2^{\lceil\log_2K\rceil}\ge K$.
Then ask for the digits in order and prune the branches that carry no alternative.
Each alternative is reached in at most $\lceil\log_2K\rceil<\log_2K+1$ distinctions.
\end{proof}

\subsection{Definition and value of the entropy}

\begin{definition}[Entropy]
\label{def:entropy}
The \emph{entropy} of the shell is
\begin{equation}
S(E,V,N)=\kB\ln K(E,V,N),
\label{eq:entropy}
\end{equation}
where $K$ is the count of Definition~\ref{def:count} and $\kB>0$ is a constant that fixes the unit of entropy.
\end{definition}

By Theorem~\ref{thm:locating}, $S/(\kB\ln2)$ is the minimum average number of binary distinctions needed to locate the state of the gas among the resolved states of its description, to within one distinction.
The equal weights required by Theorem~\ref{thm:locating} are those of Section~\ref{sec:weights}.
Each binary distinction contributes $\kB\ln2$ to the entropy.
Definition~\ref{def:entropy} therefore gives ingredient~\ref{ing:log} an operational meaning: the entropy is the cost of locating the state, in units of $\kB\ln2$ per distinction.
The count $K$ need not be an integer, and replacing it by an integer part changes $S$ by less than $\kB\ln(1+1/\lfloor K\rfloor)$.

\begin{proposition}[Entropy of the ideal gas]
\label{prop:entropy-explicit}
With $a=3N/2$ and $\delta=\Delta E/E$,
\begin{equation}
\frac{S}{\kB}=N\ln V+\ln\omega_{3N}+a\ln(2mE)+\ln\bigl[1-(1-\delta)^{a}\bigr]-\ln N!-3N\ln h.
\label{eq:entropy-explicit}
\end{equation}
\end{proposition}

\begin{proof}
Insert~\eqref{eq:shell-measure} into~\eqref{eq:count} and take the logarithm.
\end{proof}

\begin{proposition}[Role of the resolution]
\label{prop:hdep}
The resolution $h$ enters~\eqref{eq:entropy-explicit} only through the additive term $-3N\kB\ln h$.
Every derivative of $S$ with respect to $E$ or $V$ is independent of $h$.
At fixed $E$, $V$ and $N$, $S\to+\infty$ as $h\to0^{+}$.
\end{proposition}

\begin{proof}
The first two statements follow by inspection of~\eqref{eq:entropy-explicit}, and the third because $-3N\ln h\to+\infty$.
\end{proof}

Proposition~\ref{prop:hdep} is the precise sense in which a fixed resolution is needed.
Temperature and pressure (Section~\ref{sec:laws}) do not depend on $h$, but the entropy itself has a finite value only while $h$ is held fixed.

\subsection{Extensivity and the Sackur--Tetrode form}

\begin{proposition}[Sackur--Tetrode entropy]
\label{prop:sackur}
The entropy~\eqref{eq:entropy-explicit} satisfies
\begin{equation}
\frac{S}{\kB}=N\left[\ln\!\left(\frac VN\left(\frac{4\pi mE}{3Nh^{2}}\right)^{3/2}\right)+\frac52\right]+R_N,
\label{eq:sackur}
\end{equation}
with remainder
\begin{equation}
R_N=-\ln\bigl(\sqrt6\,\pi N\bigr)+\ln\bigl[1-(1-\delta)^{3N/2}\bigr]-\frac{\theta_1}{18N}-\frac{\theta_2}{12N},
\label{eq:remainder}
\end{equation}
for some $\theta_1,\theta_2\in(0,1)$.
\end{proposition}

\begin{proof}
For real $z>0$, Stirling's formula with remainder gives $\ln\Gamma(z+1)=z\ln z-z+\tfrac12\ln(2\pi z)+\theta/(12z)$ for some $\theta\in(0,1)$.
Apply it with $z=a=3N/2$ to $\ln\omega_{3N}=a\ln\pi-\ln\Gamma(a+1)$, and with $z=N$ to $\ln N!=\ln\Gamma(N+1)$.
Collecting terms, $a\ln\pi-a\ln a+a\ln(2mE)-3N\ln h=\tfrac{3N}{2}\ln\bigl(4\pi mE/(3Nh^{2})\bigr)$.
Further, $N\ln V-N\ln N=N\ln(V/N)$ and $a+N=\tfrac52N$.
The remaining terms are $-\tfrac12\ln(2\pi a)-\tfrac12\ln(2\pi N)=-\ln(\sqrt6\,\pi N)$ and the two Stirling remainders, with $12a=18N$.
\end{proof}

The bracket in~\eqref{eq:sackur} depends only on $V/N$ and $E/N$, and $R_N/N\to0$ as $N\to\infty$ at fixed $\delta$.
The entropy per particle is therefore a function of the volume and energy per particle alone, up to a term that vanishes as $N\to\infty$.
Equation~\eqref{eq:sackur} is the entropy of Sackur and Tetrode \citep{sackur1911,tetrode1912}.
Without the factor $N!$ of Lemma~\ref{lem:quotient}, the term $N\ln(V/N)$ would read $N\ln V$, and the entropy per particle would grow with $N$ at fixed $V/N$.
Section~\ref{sec:mixtures} quantifies the consequence.

%======================================================================
\section{Temperature, pressure and the equation of state}
\label{sec:laws}
%======================================================================

\subsection{Definitions}

\begin{definition}[Temperature and pressure]
\label{def:TP}
The \emph{temperature} $T$ and the \emph{pressure} $P$ of the shell are defined by
\begin{equation}
\frac1T=\left(\frac{\partial S}{\partial E}\right)_{V,N},\qquad
P=T\left(\frac{\partial S}{\partial V}\right)_{E,N},
\label{eq:TP}
\end{equation}
with the thickness $\Delta E$ held fixed in both derivatives.
\end{definition}

Equations~\eqref{eq:TP} are the standard thermodynamic definitions of temperature and pressure as conjugates of energy and volume.
The pressure of Definition~\ref{def:TP} is called the \emph{thermodynamic pressure}, to distinguish it from the mechanical pressure of Section~\ref{sec:mech}.
By~\eqref{eq:entropy-explicit}, $\partial S/\partial E>0$, so $T$ is positive and finite.

\subsection{The ideal gas law}

\begin{theorem}[Ideal gas law]
\label{thm:igl}
For every $N\ge1$, every $m>0$, every $h>0$ and every $\Delta E\in(0,E]$,
\begin{equation}
PV=N\kB T.
\label{eq:igl}
\end{equation}
\end{theorem}

\begin{proof}
By~\eqref{eq:entropy-explicit}, $V$ enters $S$ only through the term $N\kB\ln V$.
Hence $(\partial S/\partial V)_{E,N}=N\kB/V$, and~\eqref{eq:TP} gives $P=N\kB T/V$.
\end{proof}

Relation~\eqref{eq:igl} is exact, and it holds for one particle as well as for many.
Its proof uses only the factor $V^{N}$ in the measure of the shell.
That factor arises because the energy does not depend on positions and the positions range over the container independently.
The same proof applies to every region whose measure contains $V$ only through such a factor.

\begin{theorem}[Independence from the nature of the particles]
\label{thm:carrier}
Let each particle carry, besides its position $q_i\in\Lambda$, internal and kinetic coordinates $z_i$ in a measure space $(Z,\lambda)$, and let the energy be $H(x)=\sum_i\varphi(z_i)$ for a measurable function $\varphi\colon Z\to\R$ that does not involve positions.
Let the count of the shell $\{E-\Delta E<H\le E\}$ be finite and positive, with cells of fixed measure in $\Lambda\times Z$.
Then $PV=N\kB T$ holds exactly, whatever $Z$, $\lambda$ and $\varphi$ are.
\end{theorem}

\begin{proof}
The shell is $\Lambda^{N}\times B$ for a set $B\subseteq Z^{N}$ that does not depend on $V$.
Its measure is $V^{N}\lambda^{N}(B)$, and the proof of Theorem~\ref{thm:igl} applies unchanged.
\end{proof}

\begin{corollary}[Avogadro's law]
\label{cor:avogadro}
Two ideal gases in the sense of Theorem~\ref{thm:carrier} at the same pressure, volume and temperature contain the same number of particles, whatever their masses, internal structure or energy functions.
\end{corollary}

\begin{proof}
By Theorem~\ref{thm:carrier}, $N=PV/(\kB T)$ for each gas.
\end{proof}

\begin{corollary}[Boyle, Charles and Gay-Lussac]
\label{cor:classical}
At fixed $N$: at fixed $T$ the product $PV$ is constant; at fixed $P$ the volume is proportional to $T$; at fixed $V$ the pressure is proportional to $T$.
\end{corollary}

\begin{proof}
Each statement is~\eqref{eq:igl} with two of $P$, $V$ and $T$ related at fixed $N$.
\end{proof}

\subsection{Equipartition}

\begin{theorem}[Equipartition]
\label{thm:equipartition}
Let $a=3N/2$ and $\delta=\Delta E/E$.
The temperature of Definition~\ref{def:TP} satisfies
\begin{equation}
E=a\,\kB T\,(1-\eta),\qquad
\eta=\frac{\delta\,(1-\delta)^{a-1}}{1-(1-\delta)^{a}}\in[0,1).
\label{eq:equipartition}
\end{equation}
If $a\delta\ge1$, then
\begin{equation}
0\le\eta\le\frac{1}{(e-1)(a-1)}=\frac{2}{(e-1)(3N-2)},
\label{eq:eta-bound}
\end{equation}
where $e$ is the base of the natural logarithm.
\end{theorem}

\begin{proof}
By~\eqref{eq:shell-measure}, $S/\kB=c+\ln\bigl[E^{a}-(E-\Delta E)^{a}\bigr]$ with $c$ independent of $E$.
Differentiating at fixed $\Delta E$,
\[
\frac{1}{\kB T}=\frac{a\bigl[E^{a-1}-(E-\Delta E)^{a-1}\bigr]}{E^{a}-(E-\Delta E)^{a}}=\frac aE\cdot\frac{1-(1-\delta)^{a-1}}{1-(1-\delta)^{a}}.
\]
The last factor equals $1-\eta$ with $\eta$ as in~\eqref{eq:equipartition}, since $1-(1-\delta)^{a}-\delta(1-\delta)^{a-1}=1-(1-\delta)^{a-1}$.
It lies in $(0,1]$ because $a\ge\tfrac32$, which gives $\eta\in[0,1)$.
For the bound, $(1-\delta)^{a-1}\le e^{-\delta(a-1)}$ because $a-1>0$, and $1-(1-\delta)^{a}\ge1-e^{-a\delta}\ge1-e^{-1}$ when $a\delta\ge1$.
Hence $\eta\le\delta e^{-\delta(a-1)}/(1-e^{-1})$.
The function $\delta\mapsto\delta e^{-\delta(a-1)}$ has maximum $1/\bigl(e(a-1)\bigr)$ over $\delta>0$, and $e/(e-1)\cdot1/\bigl(e(a-1)\bigr)=1/\bigl((e-1)(a-1)\bigr)$.
\end{proof}

\begin{corollary}[Insensitivity to the thickness]
\label{cor:shell-insensitive}
If the thickness satisfies $\Delta E\ge2E/(3N)$, then $E=\tfrac32N\kB T\,(1-\eta)$ with $0\le\eta\le2/\bigl((e-1)(3N-2)\bigr)$.
In particular $E/(N\kB T)\to\tfrac32$ as $N\to\infty$, whatever the thickness, provided it is not smaller than $2E/(3N)$.
\end{corollary}

\begin{proof}
The condition is $a\delta\ge1$; apply Theorem~\ref{thm:equipartition}.
\end{proof}

\begin{remark}[Ball and shell]
\label{rem:ball}
If the shell is replaced by the ball $B_E$, which is the case $\delta=1$, then $\eta=0$ and $E=\tfrac32N\kB T$ holds exactly.
The ball contains energies far below that of the gas, which the flow never reaches because it conserves the energy, while the shell contains only energies within the resolved width.
The difference between the two choices is the difference between the entropies of Gibbs and of Boltzmann, discussed in \citep{gibbs1902,dunkel2014}.
By~\eqref{eq:eta-bound} the difference is of relative order $N^{-1}$.
\end{remark}

\begin{corollary}[General quadratic energy]
\label{cor:quadratic}
In the setting of Theorem~\ref{thm:carrier}, let $Z=\R^{f}$ with Lebesgue measure, and let $\varphi$ be a positive-definite quadratic form on $\R^{f}$.
Then $E=\tfrac f2N\kB T(1-\eta_f)$, where $\eta_f$ is given by~\eqref{eq:equipartition} with $a=fN/2$.
\end{corollary}

\begin{proof}
The set $\{z\in\R^{fN}:\sum_i\varphi(z_i)\le E\}$ is an ellipsoid, whose measure is a constant times $E^{fN/2}$.
The proof of Theorem~\ref{thm:equipartition} applies with $a=fN/2$.
\end{proof}

\subsection{The adiabatic law}

\begin{proposition}[Adiabatic law]
\label{prop:adiabatic}
Consider a change of $E$ and $V$ at fixed $N$ that keeps the entropy and the relative thickness $\delta$ constant.
Along such a change, $VE^{3/2}$ and $PV^{5/3}$ are constant.
\end{proposition}

\begin{proof}
At fixed $N$ and $\delta$,~\eqref{eq:entropy-explicit} reads $S/\kB=N\ln V+\tfrac{3N}{2}\ln E+c$ with $c$ constant.
Constant $S$ therefore means constant $VE^{3/2}$.
By Theorems~\ref{thm:igl} and~\ref{thm:equipartition}, $PV=N\kB T=\tfrac23E/(1-\eta)$, and $\eta$ depends only on $\delta$ and $N$.
Writing $E=C^{2/3}V^{-2/3}$ for the constant $C=VE^{3/2}$ gives $PV^{5/3}=\tfrac23C^{2/3}/(1-\eta)$, which is constant.
\end{proof}

%======================================================================
\section{Mechanical pressure}
\label{sec:mech}
%======================================================================

The pressure of Definition~\ref{def:TP} is defined from the entropy.
This section computes the pressure as a force on the walls and compares the two.

\begin{definition}[Mechanical pressure]
\label{def:mech}
Along a trajectory, a particle that reflects at a face transfers to that face the \emph{impulse} $2|p_\perp|$, where $p_\perp$ is its momentum component normal to the face.
Let $J(\tau)$ be the total impulse transferred to all six faces during the time interval $[0,\tau]$.
The \emph{mechanical pressure} of the trajectory is
\begin{equation}
\Pmech=\lim_{\tau\to\infty}\frac{J(\tau)}{\tau\cdot6L^{2}},
\label{eq:pmech}
\end{equation}
the time-averaged force per unit area of wall, whenever the limit exists.
\end{definition}

\begin{theorem}[Mechanical pressure]
\label{thm:mech}
For every trajectory of the ideal gas of Definition~\ref{def:ideal} that meets no edge or corner of $\Lambda$, the limit~\eqref{eq:pmech} exists and
\begin{equation}
\Pmech V=\tfrac23H_N(x),
\label{eq:pmech-value}
\end{equation}
where $H_N(x)$ is the energy of the trajectory.
The identity holds for every $N\ge1$.
\end{theorem}

\begin{proof}
Particles of Definition~\ref{def:ideal} do not interact, so the impulse is a sum over particles and over the three axes.
Fix a particle and the $x$-axis, and let $v_x=p_x/m$ be its velocity component along that axis.
If $v_x=0$, the particle never meets the faces normal to the $x$-axis, and its contribution is $0=p_x^{2}/(mL)\cdot\tau$.
If $v_x\ne0$, the magnitude $|v_x|$ is conserved, and the $x$-coordinate moves back and forth across $[0,L]$ with period $2L/|v_x|$.
In each period the particle meets each of the two faces once and transfers an impulse $2|p_x|$ at each meeting.
The number of meetings during $[0,\tau]$ is $\tau|v_x|/L+r$ with $|r|\le2$, so the impulse transferred to the two faces is $2|p_x|\bigl(\tau|v_x|/L+r\bigr)=\tau\,2p_x^{2}/(mL)+O(1)$.
Summing over the three axes and over the particles gives $J(\tau)=\tau\,\frac{2}{mL}\sum_i|p_i|^{2}+O(1)=\tau\,\frac{4}{L}H_N(x)+O(1)$.
Dividing by $6L^{2}\tau$ and letting $\tau\to\infty$ gives $\Pmech=\tfrac{2}{3}H_N(x)/L^{3}$.
\end{proof}

Relation~\eqref{eq:pmech-value} is the kinetic pressure of Clausius, here obtained as an exact time average rather than as an average over a distribution of velocities.
It needs neither Hypothesis~\ref{hyp:ergodic} nor any statement about the distribution of momenta.
When the particles interact rarely, the same identity holds for the time averages by the virial theorem, provided the interparticle forces contribute negligibly to the virial \citep{clausius1870}.

\begin{corollary}[Agreement of the two pressures]
\label{cor:agreement}
Let a trajectory have energy $E_0\in(E-\Delta E,E]$, so that it lies in the shell $W(E,\Delta E)$.
Its mechanical pressure and the thermodynamic pressure $P$ of the shell satisfy
\begin{equation}
(1-\delta)(1-\eta)\,P<\Pmech\le(1-\eta)\,P,
\label{eq:agreement}
\end{equation}
with $\eta$ as in Theorem~\ref{thm:equipartition}.
\end{corollary}

\begin{proof}
By Theorems~\ref{thm:igl} and~\ref{thm:equipartition}, $PV=N\kB T=\tfrac23E/(1-\eta)$.
By Theorem~\ref{thm:mech}, $\Pmech V=\tfrac23E_0$, and $E(1-\delta)<E_0\le E$.
\end{proof}

The mechanical and thermodynamic pressures therefore agree up to the factor $1-\eta$, which tends to one as $N\to\infty$ by Corollary~\ref{cor:shell-insensitive}, and up to the resolved width of the energy.

%======================================================================
\section{One particle}
\label{sec:single}
%======================================================================

\subsection{An isolated particle}

\begin{theorem}[An isolated particle]
\label{thm:single-isolated}
Let $N=1$ and let the particle have kinetic energy $\varepsilon$.
\begin{enumerate}[label=(\roman*),leftmargin=2.2em]
\item Its mechanical pressure satisfies $\Pmech V=\tfrac23\varepsilon$ exactly.
\item The thermodynamic pressure and temperature of the shell $W(\varepsilon,\Delta E;V,1)$ satisfy $PV=\kB T$ exactly.
\item That temperature satisfies $\kB T=\tfrac23\varepsilon/(1-\eta_1)$ with
\begin{equation}
1-\eta_1=\frac{1-(1-\delta)^{1/2}}{1-(1-\delta)^{3/2}},
\label{eq:eta-one}
\end{equation}
which increases from $\tfrac13$ as $\delta\to0^{+}$ to $1$ at $\delta=1$.
Hence $\tfrac23\varepsilon\le\kB T<2\varepsilon$, and the temperature of an isolated particle is fixed by its energy only up to a factor that depends on the resolved width of that energy.
\end{enumerate}
\end{theorem}

\begin{proof}
Part (i) is Theorem~\ref{thm:mech} with $N=1$, and part (ii) is Theorem~\ref{thm:igl} with $N=1$.
For (iii), set $a=\tfrac32$ in the proof of Theorem~\ref{thm:equipartition}, which gives $1/(\kB T)=\tfrac{3}{2\varepsilon}(1-\eta_1)$ with~\eqref{eq:eta-one}.
As $\delta\to0^{+}$, numerator and denominator of~\eqref{eq:eta-one} behave as $\delta/2$ and $3\delta/2$, so $1-\eta_1\to\tfrac13$.
At $\delta=1$ the ratio equals one.
The ratio is increasing on $(0,1]$, because with $s=(1-\delta)^{1/2}\in[0,1)$ it equals $(1-s)/(1-s^{3})=1/(1+s+s^{2})$, which decreases in $s$.
\end{proof}

A single isolated particle thus obeys the mechanical law $\Pmech V=\tfrac23\varepsilon$ exactly, and the thermodynamic law $PV=\kB T$ exactly.
What it lacks is a temperature fixed by its energy alone.
Theorem~\ref{thm:single-gas} shows that the other particles of a gas supply that temperature.

\subsection{One particle among many}

\begin{definition}[Kinetic pressure of one particle]
\label{def:p1}
In a gas of $N$ particles, the \emph{kinetic pressure} of particle $1$ along a trajectory is
\begin{equation}
p_1=\frac{2}{3V}\,\overline{\varepsilon_1},\qquad
\overline{\varepsilon_1}=\lim_{\tau\to\infty}\frac1\tau\int_0^{\tau}\varepsilon\bigl(p_1(t)\bigr)\,\dd t,
\label{eq:p1}
\end{equation}
whenever the time average $\overline{\varepsilon_1}$ exists.
\end{definition}

By Theorem~\ref{thm:mech} applied to a single particle, $p_1$ is the time-averaged pressure that particle $1$ exerts on the walls during every interval in which it moves freely.

\begin{theorem}[One particle in a gas]
\label{thm:single-gas}
Assume Hypothesis~\ref{hyp:ergodic}.
Let the gas have energy $E_0\in(E-\Delta E,E]$.
For almost every trajectory on $\Sigma_{E_0}$, the time-averaged kinetic energy of particle $1$ is $\overline{\varepsilon_1}=E_0/N$, and its kinetic pressure satisfies
\begin{equation}
(1-\delta)(1-\eta)\,\kB T<p_1V\le(1-\eta)\,\kB T,
\label{eq:p1-law}
\end{equation}
where $T$ is the temperature of the shell $W(E,\Delta E;V,N)$ and $\eta$ is as in Theorem~\ref{thm:equipartition}.
If $N\to\infty$ with $E/N$ and $V/N$ fixed, $\delta\to0$ and $N\delta\to\infty$, then $p_1V/(\kB T)\to1$.
\end{theorem}

\begin{proof}
By Theorem~\ref{thm:weights}, $\overline{\varepsilon_1}$ equals the $\nu_{E_0}$-average of $\varepsilon(p_1)$ for almost every trajectory.
The surface measure $\nu_{E_0}$ is invariant under permutations of the particles, so the $\nu_{E_0}$-averages of $\varepsilon(p_i)$ are equal for all $i$.
Their sum is the average of $H_N$, which equals $E_0$ on $\Sigma_{E_0}$, so each equals $E_0/N$.
Then $p_1V=\tfrac23E_0/N$, and~\eqref{eq:p1-law} follows from $\tfrac23E/N=\kB T(1-\eta)$, which is Theorem~\ref{thm:equipartition}, and $E(1-\delta)<E_0\le E$.
In the limit, $N\delta\to\infty$ gives $a\delta\ge1$ for all large $N$, so $\eta\to0$ by~\eqref{eq:eta-bound}, and $\delta\to0$ gives $1-\delta\to1$.
Both bounds in~\eqref{eq:p1-law} therefore tend to one.
\end{proof}

\begin{remark}
\label{rem:single}
The two conditions on $\delta$ in Theorem~\ref{thm:single-gas} can be met together, for instance with $\delta=N^{-1/2}$.
The law $p_1V=\kB T$ for one particle is therefore a statement about one particle among many.
The time average along the history of that one particle replaces the average over the gas, by Theorem~\ref{thm:weights}.
The remaining $N-1$ particles act as the reservoir that fixes the temperature, which an isolated particle lacks (Theorem~\ref{thm:single-isolated}).
\end{remark}

\subsection{The Maxwell--Boltzmann distribution}

\begin{proposition}[Distribution of one momentum]
\label{prop:mb}
Let $N\ge2$ and let $\rho_N$ be the density of the momentum $p_1$ of particle $1$ under the normalised surface measure $\nu_{E}/\nu_{E}(\Sigma_{E})$.
Then, with $d=3N$ and $R=\sqrt{2mE}$,
\begin{equation}
\rho_N(p)=\frac{\Gamma(d/2)}{\pi^{3/2}\,\Gamma\!\left(\frac{d-3}{2}\right)R^{3}}\left(1-\frac{|p|^{2}}{R^{2}}\right)^{\frac{d-5}{2}}\quad\text{for }|p|<R,
\label{eq:rho-n}
\end{equation}
and $\rho_N(p)=0$ otherwise.
If $N\to\infty$ with $\vartheta=2E/(3N)$ fixed, then for every $p\in\R^{3}$,
\begin{equation}
\rho_N(p)\to\bigl(2\pi m\vartheta\bigr)^{-3/2}\exp\!\left(-\frac{|p|^{2}}{2m\vartheta}\right).
\label{eq:mb}
\end{equation}
\end{proposition}

\begin{proof}
On $\Sigma_E$ the gradient of $H_N$ with respect to the momenta has magnitude $|p|/m=R/m$, which is constant.
Hence $\nu_E$ is a constant multiple of the product of Lebesgue measure on $\Lambda^{N}$ and the uniform surface measure on the sphere of radius $R$ in $\R^{d}$.
For a point uniformly distributed on that sphere, the first three coordinates have the density~\eqref{eq:rho-n} \citep{diaconis1987}.
For the limit, $R^{2}=2mE=3Nm\vartheta$, so $\bigl(1-|p|^{2}/R^{2}\bigr)^{(d-5)/2}=\bigl(1-|p|^{2}/(3Nm\vartheta)\bigr)^{(3N-5)/2}\to\exp\bigl(-|p|^{2}/(2m\vartheta)\bigr)$.
For the prefactor, $\Gamma(d/2)/\Gamma\bigl((d-3)/2\bigr)\sim(d/2)^{3/2}$ as $d\to\infty$, and $R^{3}=(dm\vartheta)^{3/2}$, so the prefactor tends to $\pi^{-3/2}(2m\vartheta)^{-3/2}$.
\end{proof}

By Theorem~\ref{thm:equipartition}, $\vartheta=\kB T(1-\eta)\to\kB T$ under the conditions of Corollary~\ref{cor:shell-insensitive}, so the limit~\eqref{eq:mb} is the Maxwell--Boltzmann distribution at temperature $T$ \citep{maxwell1860}.
By Theorem~\ref{thm:weights} it is also the distribution of the momentum of particle $1$ over time, along almost every trajectory.
For finite $N$ the density~\eqref{eq:rho-n} vanishes for $|p|\ge R$, since no particle can carry more than the whole energy.
For $N=1$ the momentum is confined to a sphere and has no density in $\R^{3}$, in agreement with Theorem~\ref{thm:single-isolated}.

%======================================================================
\section{Mixtures}
\label{sec:mixtures}
%======================================================================

\begin{definition}[Mixture]
\label{def:mixture}
A \emph{mixture} of $r$ species consists of $N_j\ge1$ particles of species $j$, $j=1,\dots,r$, with $N=\sum_jN_j$, in the container $\Lambda$.
The energy is a sum of position-independent single-particle terms that may differ between species.
Axiom~\ref{ax:labels} is replaced by invariance of every observable under permutations within each species.
The count of a symmetric region $A$ is $\mu_N(A)/\bigl(\prod_jN_j!\,h^{3N}\bigr)$.
\end{definition}

The denominator $\prod_jN_j!$ follows from Lemma~\ref{lem:quotient} applied to the group of permutations within species, whose order is $\prod_jN_j!$.

\begin{theorem}[Dalton's law]
\label{thm:dalton}
Define the \emph{partial pressure} of species $j$ as $P_j=N_j\kB T/V$.
Then the thermodynamic pressure of the mixture is $P=\sum_jP_j$.
\end{theorem}

\begin{proof}
The shell has measure $V^{N}$ times a factor independent of $V$.
As in Theorem~\ref{thm:igl}, $P=N\kB T/V=\sum_jN_j\kB T/V$.
\end{proof}

\begin{theorem}[Entropy of mixing]
\label{thm:mixing}
Let the container be divided by a partition into parts of volumes $V_j=x_jV$, with $x_j=N_j/N$, and let species $j$ be confined to part $j$, with the energy shell fixed.
Removing the partition increases the entropy by
\begin{equation}
\Delta S=-\kB\sum_{j=1}^{r}N_j\ln x_j.
\label{eq:mixing}
\end{equation}
If instead all particles belong to a single species, with $N_j$ particles initially in part $j$, removing the partition increases the entropy by
\begin{equation}
\Delta S_{\mathrm{same}}=-\kB\ln\frac{N!}{\prod_jN_j!}-\kB\sum_jN_j\ln x_j
=\frac{r-1}{2}\,\kB\ln(2\pi N)+\frac{\kB}{2}\sum_j\ln x_j+O\!\left(\frac{\kB}{N}\right)
\label{eq:mixing-same}
\end{equation}
as $N\to\infty$ at fixed $x_j$.
\end{theorem}

\begin{proof}
The momentum part of the shell is unchanged by the partition, so only the positional factor changes.
For distinct species, the positional measure with the partition is $\prod_jV_j^{N_j}$ and without it is $V^{N}$, and the symmetry factor $\prod_jN_j!$ is the same in both cases.
The ratio of counts is $V^{N}/\prod_jV_j^{N_j}=\prod_jx_j^{-N_j}$, which gives~\eqref{eq:mixing}.
For a single species, the symmetry factor is $N!$ in both cases.
With the partition, the labelled positional measure of the configurations with exactly $N_j$ particles in part $j$ is $\frac{N!}{\prod_jN_j!}\prod_jV_j^{N_j}$, since $N!/\prod_jN_j!$ counts the assignments of labelled particles to parts.
Without the partition it is $V^{N}$.
The ratio of counts, without the partition to with it, is $\frac{\prod_jN_j!}{N!}\prod_jx_j^{-N_j}$, which gives the first expression in~\eqref{eq:mixing-same}.
By Stirling's formula, $\ln\bigl(N!/\prod_jN_j!\bigr)=-\sum_jN_j\ln x_j-\tfrac{r-1}{2}\ln(2\pi N)-\tfrac12\sum_j\ln x_j+O(N^{-1})$, which gives the second.
\end{proof}

Equation~\eqref{eq:mixing} is the entropy of mixing of distinct gases, extensive in $N$.
For identical gases the increase~\eqref{eq:mixing-same} grows only logarithmically, so the entropy per particle does not change when the partition is removed.
Had the count omitted the factor $N!$, the configurations with the partition would have been counted with the labelled measure $\prod_jV_j^{N_j}$ and the increase for identical gases would have been $-\kB\sum_jN_j\ln x_j$.
An entropy increase proportional to $N$ on joining two samples of the same gas is the Gibbs mixing paradox.
The paradox is thus absent when states carry no labels, as Axiom~\ref{ax:labels} requires.

%======================================================================
\section{Free expansion}
\label{sec:expansion}
%======================================================================

\begin{proposition}[Rarity of the contracted state]
\label{prop:expansion}
Let $\Lambda_{\mathrm L}=[0,L/2]\times[0,L]^{2}$ be the left half of the container and let $A_{\mathrm L}=\{x\in\Sigma_E:\ q_i\in\Lambda_{\mathrm L}\text{ for all }i\}$.
\begin{enumerate}[label=(\roman*),leftmargin=2.2em]
\item $\nu_E(A_{\mathrm L})/\nu_E(\Sigma_E)=2^{-N}$.
\item Under Hypothesis~\ref{hyp:ergodic}, almost every trajectory spends the fraction $2^{-N}$ of its time in $A_{\mathrm L}$.
\item The entropy of the shell restricted to the left half is lower than that of the full shell by $N\kB\ln2$, which is $N$ binary distinctions.
\item If the flow sampled at a fixed time step $\tau_0>0$ is ergodic, the mean number of steps between successive entries into $A_{\mathrm L}$ is $2^{N}$.
\end{enumerate}
\end{proposition}

\begin{proof}
(i) The surface measure factorises into a positional factor, Lebesgue measure on $\Lambda^{N}$, and a momentum factor, and $\mathrm{vol}(\Lambda_{\mathrm L})^{N}/\mathrm{vol}(\Lambda)^{N}=2^{-N}$.
(ii) follows from (i) and Theorem~\ref{thm:weights}.
(iii) The restriction replaces $V^{N}$ by $(V/2)^{N}$ in~\eqref{eq:shell-measure}, which changes $S$ by $-N\kB\ln2$.
(iv) is Kac's lemma \citep{kac1947}: for an ergodic measure-preserving map of a probability space, the mean return time to a set of positive measure, starting in it, is the reciprocal of its measure.
\end{proof}

A gas released from the left half of the container is found there again only for the fraction $2^{-N}$ of the time.
For a macroscopic $N$ this fraction is too small for the event to be observed, although Poincar\'e recurrence guarantees that it occurs \citep{poincare1890}.
The entropy difference $N\kB\ln2$ is the cost of locating the gas in one half, one binary distinction per particle.

%======================================================================
\section{Discussion}
\label{sec:discussion}
%======================================================================

\subsection{What was assumed and what was derived}

Table~\ref{tab:status} lists the five ingredients of Section~\ref{sec:intro} and the source of each in the present derivation.

\begin{table}[ht]
\centering
\caption{Status of the ingredients of the standard derivation.}
\label{tab:status}
\begin{tabular}{@{}lll@{}}
\toprule
Ingredient & Standard status & Here \\
\midrule
\ref{ing:cell} cell of measure $h^{3}$ & quantum mechanics & unit of resolution (Axiom~\ref{ax:resolution}) \\
\ref{ing:nfact} division by $N!$ & indistinguishability & orbits (Lemma~\ref{lem:quotient}) \\
\ref{ing:equal} equal weights & postulate & dynamics and Axiom~\ref{ax:cells} (Section~\ref{sec:weights}) \\
\ref{ing:log} $S=\kB\ln K$ & Boltzmann's principle & locating cost (Theorem~\ref{thm:locating}) \\
\ref{ing:shell} shell thickness & convenience & positive resolution (Lemma~\ref{lem:noshell}) \\
\bottomrule
\end{tabular}
\end{table}

Two inputs remain assumptions.
The first is Hypothesis~\ref{hyp:ergodic}, whose status for interacting gases is discussed in Remark~\ref{rem:ergodic-status}.
It is used only in Theorems~\ref{thm:weights} and~\ref{thm:single-gas}, in part (ii) of Proposition~\ref{prop:expansion}, and in the reading of Proposition~\ref{prop:mb} as a time average.
The second is the form of the energy.
The equation of state needs only that the energy not depend on positions (Theorem~\ref{thm:carrier}), while the factor $\tfrac32$ of equipartition needs the quadratic kinetic energy~\eqref{eq:hamiltonian} in three dimensions (Corollary~\ref{cor:quadratic}).
The constant $\kB$ is a choice of unit and is not derived.

\subsection{The role of the resolution}

The resolution enters in three places.
It makes the energy constraint a band of positive width (Lemma~\ref{lem:noshell}).
It turns measures into counts (Definition~\ref{def:count}).
It gives the entropy a finite absolute value, which diverges as the resolution is refined without bound (Proposition~\ref{prop:hdep}).
Temperature, pressure, the equation of state, equipartition and the entropy of mixing do not depend on its value.
The resolution is therefore necessary for the entropy to be a number, and irrelevant to the laws that relate measurable quantities.

\subsection{The equation of state and its exactness}

The relation $PV=N\kB T$ holds exactly for the thermodynamic pressure, for every $N\ge1$ (Theorem~\ref{thm:igl}).
Its exactness reflects the factor $V^{N}$ in the measure of the shell, and nothing about the energy except that it does not involve positions (Theorem~\ref{thm:carrier}).
The approximations of the standard treatment enter only in the relations between energy, temperature and mechanical pressure, and they are bounded explicitly by~\eqref{eq:eta-bound} and~\eqref{eq:agreement}.

\subsection{Scope}

The derivation is classical and treats a cubic container with specular walls.
For a container of another shape, the positional factor is still $V^{N}$ and Theorems~\ref{thm:igl} and~\ref{thm:carrier} hold unchanged.
Theorem~\ref{thm:mech} uses the cube to compute the impulse on the faces, and for other shapes the virial theorem gives the same time-averaged relation \citep{clausius1870}.
Interactions that depend on positions modify the positional factor, and with it the equation of state.
For hard spheres of diameter $\sigma_0$, the excluded configurations have relative positional measure at most $\binom N2\frac{4\pi\sigma_0^{3}}{3V}$, so the positional factor tends to $V^{N}$ as $\sigma_0\to0$.
Quantum statistics enter through the occupation counts of Proposition~\ref{prop:occupation}, and they do not affect the count when $N/M$ is small.

\subsection{Conclusion}

Four axioms about an observer that resolves a bounded phase space at a fixed resolution, together with ergodicity on the energy surface, yield the ideal gas law and the laws of Avogadro, Dalton, Boyle, Charles and Gay-Lussac.
The same assumptions yield the Sackur--Tetrode entropy, the entropy of mixing, the adiabatic law and the Maxwell--Boltzmann distribution.
The five ingredients of the standard derivation appear as consequences: a unit, a band of positive width, the absence of labels, a superposition of dynamical weights, and a locating cost.
The equation of state is exact for every particle number, and the approximations are confined to equipartition and to the mechanical pressure, where they are bounded by a term of order $N^{-1}$.

\bibliographystyle{unsrtnat}
\bibliography{references}

\end{document}
